\documentclass[10pt, aspectratio=169]{beamer}
% Preámbulo modular para presentaciones Beamer (Modelación Estadística)
% Diseño y paleta institucional del Tecnológico de Monterrey

\documentclass[aspectratio=169,xcolor={dvipsnames,table}]{beamer}

\usepackage[utf8]{inputenc}
\usepackage[spanish,mexico]{babel}
\usepackage{amsmath,amssymb,amsfonts,mathtools}
\usepackage{graphicx}
\usepackage{listings}
\usepackage{booktabs}
\usepackage{tikz}
\usetikzlibrary{matrix,arrows,decorations.pathmorphing,shapes.geometric,positioning}

% Paleta Institucional Tecnológico de Monterrey
\definecolor{TecAzul}{HTML}{0039A6}       % Azul TEC: color primario institucional
\definecolor{TecAzulOscuro}{HTML}{1A2E51} % Azul marino oscuro
\definecolor{TecRojo}{HTML}{EC2661}       % Rojo de acento (paleta miTec)
\definecolor{TecGris}{HTML}{646464}       % Gris neutro para comentarios y subtítulos
\definecolor{TecFondoBloque}{HTML}{F4F6F9}% Fondo suave para bloques

% Configuración de Tema Beamer
\usetheme{Boadilla}
\usecolortheme[named=TecAzulOscuro]{structure}
\setbeamercolor{alerted text}{fg=TecRojo}
\setbeamercolor{palette primary}{bg=TecAzulOscuro,fg=white}
\setbeamercolor{palette secondary}{bg=TecAzul,fg=white}
\setbeamercolor{palette tertiary}{bg=TecAzulOscuro!80!black,fg=white}
\setbeamercolor{title}{fg=TecAzulOscuro,bg=TecFondoBloque}
\setbeamercolor{frametitle}{fg=TecAzulOscuro,bg=TecFondoBloque}

% Bloques personalizados elegantes
\setbeamercolor{block title}{bg=TecAzulOscuro,fg=white}
\setbeamercolor{block body}{bg=TecFondoBloque,fg=black}
\setbeamercolor{block title alerted}{bg=TecRojo,fg=white}
\setbeamercolor{block body alerted}{bg=TecRojo!8,fg=black}
\setbeamercolor{block title example}{bg=TecAzul,fg=white}
\setbeamercolor{block body example}{bg=TecAzul!8,fg=black}

% Redondear bordes y añadir sombra sutil
\setbeamertemplate{blocks}[rounded][shadow=true]
\setbeamertemplate{navigation symbols}{}

% Pie de página limpio con número de diapositiva
\setbeamertemplate{footline}{
  \leavevmode%
  \hbox{%
  \begin{beamercolorbox}[wd=.7\paperwidth,ht=2.25ex,dp=1ex,left]{author in head/foot}%
    \hspace*{2ex}\insertshortauthor~~(\insertshortinstitute) \hfill \insertshorttitle
  \end{beamercolorbox}%
  \begin{beamercolorbox}[wd=.3\paperwidth,ht=2.25ex,dp=1ex,right]{date in head/foot}%
    \insertshortdate{}\hspace*{2em}
    \insertframenumber{} / \inserttotalframenumber\hspace*{2ex}
  \end{beamercolorbox}}%
  \vskip0pt%
}

% Configuración de Listings de Python para visualización pedagógica de código
\lstset{
    language=Python,
    basicstyle=\ttfamily\footnotesize,
    keywordstyle=\color{TecAzulOscuro}\bfseries,
    stringstyle=\color{TecRojo},
    commentstyle=\color{TecGris}\itshape,
    showstringspaces=false,
    breaklines=true,
    frame=lines,
    rulecolor=\color{TecAzulOscuro!30},
    backgroundcolor=\color{TecFondoBloque},
    aboveskip=0.5em,
    belowskip=0.5em,
    literate={á}{{\'a}}1 {é}{{\'e}}1 {í}{{\'i}}1 {ó}{{\'o}}1 {ú}{{\'u}}1
             {Á}{{\'A}}1 {É}{{\'E}}1 {Í}{{\'I}}1 {Ó}{{\'O}}1 {Ú}{{\'U}}1
             {ñ}{{\~n}}1 {Ñ}{{\~N}}1 {¿}{{?`}}1 {¡}{{!`}}1
}

% Metadatos institucionales por defecto
\title[Modelación Estadística]{Modelación Estadística para Ciencia de Datos e Ingeniería}
\author[J. Castillo Colmenares]{Juliho David Castillo Colmenares}
\institute[Tec de Monterrey]{Tecnológico de Monterrey \\ Escuela de Ingeniería y Ciencias}
\date{\today}

% Comandos y macros estadísticas compartidas para las presentaciones Beamer (Metropolis)

% Conjuntos numéricos básicos
\newcommand{\R}{\mathbb{R}}
\newcommand{\N}{\mathbb{N}}
\newcommand{\Q}{\mathbb{Q}}
\newcommand{\Z}{\mathbb{Z}}
\newcommand{\set}[1]{\left\{ #1 \right\}}
\newcommand{\abs}[1]{\left|#1\right|}
\newcommand{\norm}[1]{\left\| #1 \right\|}

% Comandos estadísticos y probabilísticos del curso
\newcommand{\Var}{\operatorname{Var}}
\newcommand{\cov}{\operatorname{Cov}}
\newcommand{\corr}{\rho}
\newcommand{\comb}[2]{\begin{pmatrix} #1 \\ #2 \end{pmatrix}}
\newcommand{\s}{\sigma}
\newcommand{\particion}{\mathfrak{P}}
\newcommand{\card}[1]{\#\left( #1 \right)}
\newcommand{\seq}[3]{\set{#1}_{#2}^{#3}}
\newcommand{\E}{\mathbb{E}}
\newcommand{\Prob}{\mathbb{P}}

% Entornos pedagógicos y cajas en español académico natural
\newenvironment{discusion}{%
  \begin{alertblock}{Pregunta de Discusión}%
}{%
  \end{alertblock}%
}

\newenvironment{reflexion}{%
  \begin{alertblock}{Reflexión para el Aula}%
}{%
  \end{alertblock}%
}

\newenvironment{paradoja}{%
  \begin{alertblock}{Paradoja de Exploración}%
}{%
  \end{alertblock}%
}

\newenvironment{motivacion}{%
  \begin{exampleblock}{Motivación: Ciencia de Datos en la Práctica}%
}{%
  \end{exampleblock}%
}

\newenvironment{retoclase}{%
  \begin{block}{Reto Rápido en Equipo}%
}{%
  \end{block}%
}


\title[Diagnóstico de Regresión]{Sección 09.09: Diagnóstico de Residuos y Multicolinealidad}
\subtitle{Supuestos Clásicos, VIF, y Observaciones Influyentes}
\author[J. Castillo Colmenares]{Juliho Castillo Colmenares}
\institute{Tecnológico de Monterrey}
\date{\vspace{-1.2cm}}

\begin{document}

% Slide 1: Portada
\begin{frame}[plain]
  \titlepage
\end{frame}

% Slide 2: Hoja de Ruta
\begin{frame}{Hoja de Ruta --- Capítulo 09: Regresiones Lineales y Múltiples}
  \scriptsize
  ¿Dónde nos encontramos en el desarrollo modular del capítulo?
  \begin{itemize}\itemsep=0.02em
    \item \textbf{09.01} Correlación como Premisa de la Regresión
    \item \textbf{09.02} Introducción a la Regresión Lineal
    \item \textbf{09.03} Matemáticas de la Regresión: Derivación MCO y $R^2$
    \item \textbf{09.04} Regresión Lineal sobre Datos Simulados
    \item \textbf{09.05} Coeficientes Óptimos, Pruebas $t$/$F$ y RSE
    \item \textbf{09.06} Implementación en Python con \texttt{statsmodels}
    \item \textbf{09.07} Regresión Múltiple, Ecuación Normal y Ridge/Lasso
    \item \textbf{09.08} Validación de Modelos y $k$-fold Cross-Validation
    \item \textbf{09.09} Diagnóstico de Residuos y Multicolinealidad (VIF) \emph{(Hoy)}
    \item \textbf{09.10} Regresión con \texttt{scikit-learn}
    \item \textbf{09.11} Variables Categóricas y Variables Muda
    \item \textbf{09.12} Transformaciones No Lineales y Regresión Polinomial
  \end{itemize}
  \pause
  \vspace{-0.05cm}
  \begin{block}{Objetivo de la Sesión}
    Auditar rigurosamente si un modelo de regresión ajustado es \emph{confiable}: verificar los supuestos clásicos mediante análisis de residuos y pruebas formales (Durbin-Watson, Breusch-Pagan), cuantificar la multicolinealidad con el Factor de Inflación de Varianza, e identificar observaciones influyentes mediante la Distancia de Cook.
  \end{block}
\end{frame}

% Slide 3: Motivación
\begin{frame}{¿Es Confiable el Modelo que Ajustamos?}
  \footnotesize
  \begin{motivacion}
    En las Secciones 09.03 a 09.07 construimos y resolvimos el modelo de regresión (simple y múltiple), pero toda inferencia (intervalos, pruebas $t$/$F$) descansa en supuestos sobre el término de error $\varepsilon$ que \textbf{nunca se verifican automáticamente}. Un modelo con $R^2$ alto puede ser completamente inválido si sus residuos violan estos supuestos.
  \end{motivacion}

  \vspace{0.15cm}
  \pause
  \begin{itemize}\itemsep=0.1cm
    \item Los \textbf{residuos} $e_i=Y_i-\hat Y_i$ son la principal herramienta diagnóstica. \pause
    \item Complementamos el análisis gráfico con pruebas formales y con la identificación de observaciones problemáticas.
  \end{itemize}
\end{frame}

% Slide 4: Supuestos clásicos y residuos
\begin{frame}{Los Supuestos Clásicos y el Análisis de Residuos}
  \footnotesize
  \begin{block}{Cinco Supuestos sobre $\varepsilon_i$}
    Linealidad, Independencia, Homocedasticidad ($\text{Var}(\varepsilon_i)=\sigma^2$), Normalidad ($\varepsilon_i\sim N(0,\sigma^2)$), y Media Cero.
  \end{block}
  \pause

  \vspace{0.1cm}
  \begin{alertblock}{Patrones en la Gráfica de Residuos vs. Ajustados}
    Forma de U $\to$ no linealidad; \textbf{forma de embudo} $\to$ heterocedasticidad; puntos aislados extremos $\to$ posibles atípicos.
  \end{alertblock}
\end{frame}

% Slide 5: Pruebas formales
\begin{frame}{Pruebas Formales: Durbin-Watson y Breusch-Pagan}
  \footnotesize
  \begin{block}{Durbin-Watson (Independencia)}
    $$DW=\frac{\sum_{i=2}^n(e_i-e_{i-1})^2}{\sum_{i=1}^n e_i^2}\approx2(1-\hat\rho), \qquad DW\in(0,4),\ DW\approx2 \text{ ideal}$$
  \end{block}
  \pause

  \vspace{0.1cm}
  \begin{alertblock}{Breusch-Pagan (Homocedasticidad)}
    Regresión auxiliar de $\hat\varepsilon_i^2$ sobre los predictores: $$LM=n\cdot R_{\text{aux}}^2 \sim \chi^2_p \text{ bajo } H_0\text{: homocedasticidad}$$
  \end{alertblock}
\end{frame}

% Slide 6: VIF
\begin{frame}{Multicolinealidad: el Factor de Inflación de Varianza}
  \footnotesize
  \begin{block}{Definición Operativa}
    $$\text{VIF}_j=\frac{1}{1-R_j^2}, \qquad R_j^2: \text{ coef. de regredir } X_j \text{ contra los demás predictores}$$
  \end{block}
  \pause

  \vspace{0.1cm}
  \begin{alertblock}{Identidad con la Matriz de Correlación}
    $\text{VIF}_j=(\mathbf{R}^{-1})_{jj}$: un VIF alto es la manifestación diagnóstica de que $\mathbf{X}^T\mathbf{X}$ se aproxima a la singularidad (Sección 09.07). Umbral convencional: $\text{VIF}>5$ amerita Ridge, Lasso, o eliminación de variables.
  \end{alertblock}
\end{frame}

% Slide 7: Observaciones influyentes
\begin{frame}{Observaciones Influyentes: la Distancia de Cook}
  \footnotesize
  \begin{block}{Fórmula Computacional}
    $$D_i=\frac{r_i^2}{p+1}\cdot\frac{h_{ii}}{1-h_{ii}}, \qquad \text{umbral: } D_i>4/n$$
  \end{block}
  \pause

  \vspace{0.1cm}
  \begin{alertblock}{Dos Mecanismos Distintos}
    Alto apalancamiento $h_{ii}$ (posición extrema en $X$) \textbf{o} residuo grande $r_i$ (posible atípico en $Y$) --- cualquiera de los dos, o ambos, puede disparar $D_i$.
  \end{alertblock}
\end{frame}

% Slide 8: Lab Python (Bloque 1)
\begin{frame}[fragile]{Laboratorio en Python: Durbin-Watson y Breusch-Pagan}
  \scriptsize
  Ambas pruebas implementadas desde cero con \texttt{numpy}/\texttt{scipy} (\texttt{09.03\_regression\_diagnostics.py}):
  {\lstset{basicstyle=\fontsize{3.3pt}{4.0pt}\selectfont\ttfamily,aboveskip=0.05em,belowskip=0.05em}
  \lstinputlisting[firstline=19, lastline=41]{../../code/09_regresiones/09.09_regression_diagnostics.py}}
\end{frame}

% Slide 9: Lab Python (Bloque 2)
\begin{frame}[fragile]{Laboratorio en Python: VIF y su Identidad Matricial}
  \scriptsize
  VIF calculado por regresión auxiliar y verificado contra $\text{diag}(\mathbf{R}^{-1})$:
  {\lstset{basicstyle=\fontsize{3.4pt}{4.1pt}\selectfont\ttfamily,aboveskip=0.05em,belowskip=0.05em}
  \lstinputlisting[firstline=44, lastline=70]{../../code/09_regresiones/09.09_regression_diagnostics.py}}
\end{frame}

% Slide 10: Lab Python (Bloque 3)
\begin{frame}[fragile]{Laboratorio en Python: Distancia de Cook}
  \scriptsize
  Detección de observaciones influyentes por alto apalancamiento vs. residuo grande:
  {\lstset{basicstyle=\fontsize{3.2pt}{3.9pt}\selectfont\ttfamily,aboveskip=0.05em,belowskip=0.05em}
  \lstinputlisting[firstline=73, lastline=103]{../../code/09_regresiones/09.09_regression_diagnostics.py}}
\end{frame}

% Slide 11: Salida Terminal
\begin{frame}[fragile]{Laboratorio en Python: Salida en Terminal}
  \scriptsize
  Salida estándar generada al ejecutar el script del laboratorio en Python:
  \vspace{0.02cm}
  \begin{block}{Salida Estándar en Consola}
    \fontsize{3.6pt}{4.3pt}\selectfont
    \begin{verbatim}
=== Block 1: Durbin-Watson and Breusch-Pagan ===
DW=0.5759 (approx 2(1-rho)=0.7866) -> strong positive autocorrelation
Breusch-Pagan: R2_aux=0.1760, LM=10.5576, p=0.0051 -> heteroscedasticity

=== Block 2: VIF ===
Predictor 1: VIF=94.59  Predictor 2: VIF=94.55  Predictor 3: VIF=1.01
VIF via 1/(1-R^2) matches VIF via diag(R^-1) exactly

=== Block 3: Cook's Distance ===
Obs 0: leverage=0.59, resid=-0.85, D=0.52 -> INFLUENTIAL (leverage-driven)
Obs 1: leverage=0.04, resid=5.10,  D=0.54 -> INFLUENTIAL (residual-driven)
    \end{verbatim}
  \end{block}
\end{frame}

% Slide 12A: Ejercicio Nivel 1 (Enunciado)

% Slide 12B: Ejercicio Nivel 1 (Resolución)

% Slide 13A: Ejercicio Nivel 2 (Enunciado)

% Slide 13B: Ejercicio Nivel 2 (Resolución)

% Slide 14A: Ejercicio Nivel 3 (Enunciado)
\begin{frame}{Ejercicio en Clase --- Nivel 3: Analítico (1/2)}
  \footnotesize
  \begin{block}{Problema --- Distancia de Cook (Enunciado, Problema 9.14.7)}
    Con $n=50$, $p=4$ ($p+1=5$): Obs. A ($h_{ii}=0.15,r_i=1.2$), Obs. B ($h_{ii}=0.45,r_i=0.8$), Obs. C ($h_{ii}=0.08,r_i=3.5$). Calcule $D_i$ para cada una y decida cuáles son influyentes ($4/n=0.08$).
  \end{block}

  \vspace{0.1cm}
  \pause
  \begin{alertblock}{Estrategia}
    $D_i=\dfrac{r_i^2}{p+1}\cdot\dfrac{h_{ii}}{1-h_{ii}}$ para cada observación.
  \end{alertblock}
\end{frame}

% Slide 14B: Ejercicio Nivel 3 (Resolución)
\begin{frame}{Ejercicio en Clase --- Nivel 3: Analítico (2/2)}
  \scriptsize
  Sustituimos en la fórmula computacional: \pause

  \vspace{0.05cm}
  \begin{block}{Cálculo}
    \scriptsize
    $$D_A=\frac{1.44}{5}(0.1765)\approx0.051, \quad D_B=\frac{0.64}{5}(0.8182)\approx0.105, \quad D_C=\frac{12.25}{5}(0.0870)\approx0.213$$
  \end{block}

  \vspace{0.05cm}
  \pause
  \begin{alertblock}{Interpretación}
    \scriptsize
    $D_A<0.08$ (no influyente); $D_B,D_C>0.08$ (influyentes). B lo es por \textbf{alto apalancamiento} pese a residuo moderado; C lo es por \textbf{residuo grande} pese a apalancamiento bajo: dos mecanismos distintos de influencia.
  \end{alertblock}
\end{frame}

% Slide 15A: Ejercicio Nivel 4 (Enunciado)
\begin{frame}{Ejercicio en Clase --- Nivel 4: Desafiante (1/2)}
  \footnotesize
  \begin{block}{Problema --- Deducción de $DW\approx2(1-\hat\rho)$ (Enunciado, Problema 9.14.9)}
    Demuestre que $DW\approx2(1-\hat\rho)$ donde $\hat\rho=\sum e_ie_{i-1}/\sum e_i^2$, y deduzca el rango teórico $DW\in(0,4)$ con sus tres casos límite.
  \end{block}

  \vspace{0.1cm}
  \pause
  \begin{alertblock}{Estrategia}
    \scriptsize
    Expanda $(e_i-e_{i-1})^2$ y aproxime las sumas parciales por la suma completa $\sum e_i^2$.
  \end{alertblock}
\end{frame}

% Slide 15B: Ejercicio Nivel 4 (Resolución)
\begin{frame}{Ejercicio en Clase --- Nivel 4: Desafiante (2/2)}
  \scriptsize
  Expandimos el numerador y sustituimos: \pause

  \vspace{0.05cm}
  \begin{block}{Cálculo}
    \scriptsize
    $$\sum(e_i-e_{i-1})^2\approx2\sum e_i^2-2\sum e_ie_{i-1} \implies DW\approx2-2\hat\rho=2(1-\hat\rho)$$
  \end{block}

  \vspace{0.05cm}
  \pause
  \begin{alertblock}{Interpretación}
    \scriptsize
    Como $\hat\rho\in[-1,1]$: $\hat\rho\to1\implies DW\to0$ (autocorrelación positiva perfecta); $\hat\rho=0\implies DW=2$ (ideal); $\hat\rho\to-1\implies DW\to4$ (autocorrelación negativa perfecta).
  \end{alertblock}
\end{frame}

% Slide 16: Cierre
\begin{frame}{Cierre de Sección: Diagnóstico de Regresión}
  \begin{reflexion}
    Un modelo ajustado nunca es confiable por decreto: hemos auditado los supuestos clásicos mediante residuos y pruebas formales (Durbin-Watson, Breusch-Pagan), cuantificado la multicolinealidad con el VIF, y detectado observaciones influyentes mediante la Distancia de Cook --- cerrando el ciclo completo de validación interna del modelo.
  \end{reflexion}

  \vspace{0.2cm}
  \pause
  \begin{block}{Lecciones Clave}
    \begin{itemize}\itemsep=0.08cm
      \item \textbf{$DW\approx2$:} valor ideal; desviaciones indican autocorrelación positiva o negativa.
      \item \textbf{VIF$_j=(\mathbf{R}^{-1})_{jj}$:} un VIF alto refleja la misma inestabilidad numérica de $\mathbf{X}^T\mathbf{X}$ vista en la Sección 09.07.
      \item \textbf{Distancia de Cook:} apalancamiento alto y residuo grande son mecanismos de influencia independientes.
    \end{itemize}
  \end{block}
\end{frame}

% Slide 17: Perspectiva Modular
\begin{frame}{Perspectiva Modular --- El Próximo Reto}
  \scriptsize
  \begin{retoclase}
    Con el modelo auditado internamente y ya validado externamente en la Sección 09.08, cerramos el núcleo estadístico del capítulo. Las tres secciones finales (09.10-09.12) son de naturaleza más aplicada: la Sección 09.10 muestra cómo todo lo construido a mano en \texttt{numpy}/\texttt{scipy} se traduce a la práctica profesional con \texttt{scikit-learn}.
  \end{retoclase}

  \vspace{0.08cm}
  \pause
  \begin{alertblock}{Preparación Recomendada}
    \begin{itemize}\itemsep=0.06cm
      \item Reflexiona sobre por qué un modelo puede pasar todos los diagnósticos de esta sección y aun así generalizar mal (sobreajuste) --- por eso la validación de la Sección 09.08 es independiente de este diagnóstico interno.
      \item Repasa la API básica de \texttt{scikit-learn} (\texttt{fit}, \texttt{predict}, \texttt{score}) que unifica la interfaz de todos los modelos de la librería.
    \end{itemize}
  \end{alertblock}
\end{frame}

\end{document}
